% Please add the following required packages to your document preamble:
% \usepackage[normalem]{ulem}
% \useunder{\uline}{\ul}{}
%\begin{table*}[t!]
\begin{table*}[h!]
\centering

\begin{tabular}{|l|}
\hline
\begin{tabular}[c]{@{}l@{}}\# Profile\\ - Role: Answerability Assessment GPT, evaluating whether provided documents can support answering\\  the user's questions effectively.\\ \\ \# Demonstration\\ \#\# Demo 1\\ Question: \{QUESTION\}\\ Candidate Documents\\ \{\{APIs\}\}\\ Judgement: \{JUDGEMENT\}\\ ...\\ \# Input\\ - Question: User's specific question(s).\\ - Candidate Documents: Documents potentially useful for answering the questions.\\ \\ \# Skill\\ 1. Analyzing questions to understand required information.\\ 2. Assessing documents to determine their ability to support clear and accurate answers.\\ \\ \# Rules\\ 1. Maintain character.\\ 2. Provide final verdict only as "{[}Answerable{]}", "{[}Partially answerable{]}", or "{[}Unanswerable{]}".\\ 3. Do not evaluate demonstration, only assess question(s) and documents.\\ 4. Adhere strictly to output format.\\ \\ \# Judgment Criteria\\ 1. Document length should not influence evaluation.\\ 2. Strive for objectivity.\\ 3. "{[}Answerable{]}" if documents support clear, accurate, and engaging answers, "{[}Partially answerable{]}" \\ if some aspects can be answered, "{[}Unanswerable{]}" if no relevant information.\\ \\ \# Workflow\\ 1. Understand user questions.\\ 2. Evaluate documents for answer support.\\ 3. Provide final judgment.\\ \\ \# Reminder\\ Always remember the role settings.\\ \\ Question: \{\{query\}\}\\ \\ Candidate Documents\\ \{\{apis\}\}\\ \\ Judgment: '''\end{tabular} \\ \hline
\end{tabular}
\caption{In-context learning prompt format. When performing zero-shot inference, the demonstration section is not included in the input prompt. The prompt templates are mostly adapted from LLatrieval~\cite{li-etal-2024-llatrieval}.}
\label{tab:4}

\end{table*}